\begin{lemma}[Depth-integral slab-volume bound]\label{lem:rolling}
Let $K$ have the interior rolling ball property at radius $\rho_K$, and let
$n\ge2$.  Then for every $x\in K$, every unit vector $u$, and every
$0<L\le\rho_K$,
\[
\vol\Bigl(K\cap\bigl\{y:\ |\langle y-x,u\rangle|\le L\bigr\}\Bigr)
\;\ge\;
a_n\,\rho_K^{(n-1)/2}\,L^{(n+1)/2},
\qquad
a_n:=\frac{\omega_{n-1}}{2^{\,n}},
\]
where $\omega_{n-1}$ is the volume of the unit ball of $\R^{n-1}$.  (For $n=1$
the left-hand side is at least $L$, and the statement holds trivially with
$a_1=1$.)
\end{lemma}

\begin{proof}
Write $\rho:=\rho_K$, $\tau(y):=\langle y,u\rangle$,
$\tau_{\max}:=\max_K\tau$, $\tau_{\min}:=\min_K\tau$, and let
$A(\tau):=\vol_{n-1}\bigl(K\cap\{\tau(y)=\tau\}\bigr)$, so that the slab volume
equals $\int_{\tau(x)-L}^{\tau(x)+L}A(\tau)\,d\tau$.  Note that $K$ contains a
ball of radius $\rho$, so $\tau_{\max}-\tau_{\min}\ge2\rho\ge2L$.  By the
Brunn--Minkowski inequality, $A^{1/(n-1)}$ is concave on
$[\tau_{\min},\tau_{\max}]$; in particular $A$ is \emph{unimodal}, so on any
subinterval $A$ is bounded below by the smaller of its endpoint values.

\emph{Step 1: an anchor level near each extreme.}  Let $p\in\partial K$ attain
$\tau_{\max}$ and let $B=B(c,\rho)\subseteq K$ be a rolling ball with
$p\in\partial B$; set $g:=\tau(p)-\tau(c)\in[0,\rho]$ (nonnegative since
$\tau(c)\le\tau_{\max}$, and at most $|p-c|=\rho$).  The section of $B$ at the
level $\tau_{\max}-h$ is an $(n-1)$-ball of radius $r(h)$ with
$r(h)^2=\rho^2-(g-h)^2$.  For $0\le h\le g$,
\[
r(h)^2=(\rho-g+h)(\rho+g-h)\;\ge\;h\cdot\rho,
\]
since $\rho-g+h\ge h$ and $\rho+g-h\ge\rho$.  Hence
\begin{equation}\label{eq:section}
A(\tau_{\max}-h)\;\ge\;\omega_{n-1}(\rho h)^{(n-1)/2}
\qquad(0\le h\le g),
\end{equation}
and at $h=g$ the section of $B$ is its full equator:
$A(\tau_{\max}-g)\ge\omega_{n-1}\rho^{\,n-1}$.

Define the \emph{anchor depth} $a:=\min(g,\,L/4)$ and the value
$A^*:=\omega_{n-1}(\rho L/4)^{(n-1)/2}$.  Then
\begin{equation}\label{eq:anchor}
A(\tau_{\max}-a)\;\ge\;A^* :
\end{equation}
if $g\ge L/4$ then $a=L/4\le g$ and \eqref{eq:section} gives
$A(\tau_{\max}-a)\ge\omega_{n-1}(\rho L/4)^{(n-1)/2}=A^*$; if $g<L/4$ then
$a=g$ and the equator bound gives
$A(\tau_{\max}-a)\ge\omega_{n-1}\rho^{n-1}\ge A^*$, using $\rho\ge L/4$.
Symmetrically, with $g'$ the offset of a rolling ball at a minimizer of
$\tau$, the anchor depth $a':=\min(g',L/4)$ satisfies
$A(\tau_{\min}+a')\ge A^*$.  Note $a,a'\le L/4$.

\emph{Step 2: the integral.}  By \eqref{eq:anchor} and unimodality,
\[
A(\tau)\;\ge\;A^*
\qquad\text{for all }
\tau\in M:=[\tau_{\min}+a',\ \tau_{\max}-a].
\]
The complement of $M$ in $[\tau_{\min},\tau_{\max}]$ has total length
$a+a'\le L/2$.  The slab's level interval
$[\tau(x)-L,\tau(x)+L]$ intersects $[\tau_{\min},\tau_{\max}]$ in an interval
of length at least $L$: indeed $\tau(x)$ lies in the range, and the range has
length at least $2\rho\ge2L$, so at least one of the two sides of $\tau(x)$
contains a full segment of length $L$ inside the range.  Consequently the slab
meets $M$ in a set of length at least $L-\frac L2=\frac L2$, and
\[
\int_{\tau(x)-L}^{\tau(x)+L}A
\;\ge\;\frac L2\,A^*
\;=\;\frac{\omega_{n-1}}{2\cdot4^{(n-1)/2}}\,\rho^{(n-1)/2}L^{(n+1)/2}
\;=\;\frac{\omega_{n-1}}{2^{\,n}}\,\rho^{(n-1)/2}L^{(n+1)/2}.
\]
\end{proof}
